\documentclass[12pt]{article}
\usepackage[utf8]{inputenc}
\usepackage{amsmath,amssymb,amsthm}
\usepackage{algorithm}
\usepackage{algpseudocode}
\usepackage{geometry}
\geometry{margin=1in}

\newtheorem{theorem}{Theorem}
\newtheorem{lemma}[theorem]{Lemma}
\newtheorem{corollary}[theorem]{Corollary}

\title{Problem 111: Primes with Runs}
\author{Project Euler Solution}
\date{}

\begin{document}
\maketitle

\section{Problem Statement}

For $n$-digit primes, define:
\begin{itemize}
    \item $M(n, d)$: the maximum number of times digit $d$ can appear in an $n$-digit prime,
    \item $N(n, d)$: the count of $n$-digit primes achieving $M(n, d)$ repetitions of $d$,
    \item $S(n, d)$: the sum of all such primes.
\end{itemize}
Compute $\displaystyle S(10) = \sum_{d=0}^{9} S(10, d)$.

\section{Mathematical Development}

\begin{theorem}[Repdigit composite bound]
A repdigit number $\underbrace{dd \cdots d}_{n}$ with $n \ge 2$ and $1 \le d \le 9$ is composite.
\end{theorem}

\begin{proof}
Write $\underbrace{dd \cdots d}_{n} = d \cdot R_n$ where $R_n = (10^n - 1)/9$ is the repunit of length~$n$. For $n \ge 2$ we have $R_n \ge 11$, so the factorization $d \cdot R_n$ is nontrivial.
\end{proof}

\begin{corollary}
For all $n \ge 2$ and digits $d \in \{0,1,\ldots,9\}$, $M(n,d) \le n - 1$.
\end{corollary}

\begin{proof}
If $d = 0$, an $n$-digit number cannot have all $n$ digits equal to~0. If $d \ge 1$, the repdigit $\underbrace{dd\cdots d}_n$ is composite by Theorem~1. In both cases, no $n$-digit prime consists of a single repeated digit.
\end{proof}

\begin{theorem}[Candidate enumeration]
For fixed $n$, $d$, and $k$ (the number of positions holding digit~$d$), the raw candidate count is $\binom{n}{n-k} \cdot 9^{n-k}$, from which those with a leading zero are discarded.
\end{theorem}

\begin{proof}
Select which $n - k$ of the $n$ positions hold a digit other than~$d$: $\binom{n}{n-k}$ ways. Each such position admits 9 choices from $\{0,\ldots,9\}\setminus\{d\}$, yielding $\binom{n}{n-k}\cdot 9^{n-k}$ candidates total. Candidates with leading digit~0 are invalid $n$-digit numbers and are excluded.
\end{proof}

\begin{lemma}[Deterministic Miller--Rabin]
The Miller--Rabin test with witnesses $\{2,3,5,7,11,13\}$ is deterministic for all integers below $3.2 \times 10^{13}$ (Jaeschke, 1993). This covers every 10-digit number.
\end{lemma}

\begin{theorem}[Greedy descent correctness]
Starting from $k = n-1$ and descending, the first $k$ at which at least one $n$-digit prime with exactly $k$ copies of digit~$d$ is found equals $M(n,d)$.
\end{theorem}

\begin{proof}
By Corollary~1, $M(n,d) \le n-1$. At each level $k$, the algorithm exhaustively enumerates all valid candidates (Theorem~2) and applies a correct primality test (Lemma~3). Therefore the first non-empty level is the maximum, which is $M(n,d)$ by definition.
\end{proof}

\section{Algorithm}

\begin{algorithm}[H]
\caption{Compute $S(10)$}
\begin{algorithmic}[1]
\State $\text{total} \gets 0$
\For{$d = 0$ to $9$}
    \For{$k = 9$ \textbf{downto} $1$}
        \State $\text{primes} \gets \emptyset$
        \ForAll{subsets $P \subseteq \{0,\ldots,9\}$ with $|P| = 10-k$}
            \ForAll{digit assignments to positions in $P$ from $\{0,\ldots,9\}\setminus\{d\}$}
                \State Construct number $N$ with digit $d$ at all positions $\notin P$
                \If{leading digit of $N$ is nonzero \textbf{and} $\textsc{IsPrime}(N)$}
                    \State $\text{primes} \gets \text{primes} \cup \{N\}$
                \EndIf
            \EndFor
        \EndFor
        \If{$\text{primes} \neq \emptyset$}
            \State $\text{total} \gets \text{total} + \sum_{p \in \text{primes}} p$
            \State \textbf{break}
        \EndIf
    \EndFor
\EndFor
\State \Return total
\end{algorithmic}
\end{algorithm}

\section{Complexity Analysis}

\begin{itemize}
\item \textbf{Time:} At worst $k = 8$ (two free positions), giving $\binom{10}{2} \cdot 9^2 = 3645$ candidates per digit and at most $36{,}450$ total. Each Miller--Rabin test with 6 witnesses costs $O(\log^2 n)$ where $n < 10^{10}$, which is bounded by a constant. Overall: $O(1)$.
\item \textbf{Space:} $O(1)$ auxiliary storage.
\end{itemize}

\section{Answer}

\[
\boxed{612407567715}
\]

\end{document}
